\begin{abstract}
What can a synthetic fraud benchmark show, and what can it not? We answer the question for
FraudShield-EAC, a seeded synthetic benchmark of East African digital payments (mobile money,
USSD, agent banking, card, online and bank transfer, across five countries) built to develop and
evaluate a fraud detection system. Its clearest negative result is a pre-registered variant based
on a scam documented among Kenyan mobile-money users, modelled as a single victim-initiated
transfer to a counterparty the victim already knows, so that it has neither burst structure nor
counterparty novelty: a diagnostic XGBoost model catches it at \nRevRecall{} \nRevCI{}, against \nBaseVarRecall{}
\nBaseVarCI{} for the base scenarios.

The reason is the benchmark's structure. A single engineered feature, the one-hour to thirty-day
velocity ratio, separates the classes at \nFloorM{} on its own, so we report every model metric as
a margin over that single-feature floor rather than over chance. Trained on \nTrainUsed{} rows, the
calibrated XGBoost and LightGBM ensemble reaches an AUC of \nGateAUC{} \nGateAUCci{}, a margin of
\nGateMargin{} over the floor on the same test rows. Four disjoint feature groups each reach at
least \nGroupGeographic{} alone, so ablations on this benchmark measure redundancy, not importance,
and both generalisation experiments we planned (leave-one-country-out and a novel fraud variant
held out in time) are null: the generator encodes fraud as bursts, and every country and variant is
a view of that structure. One result does not depend on how easy the data is: the median cost of
an exact TreeSHAP explanation grows with tree complexity, about \nFrontRatio-fold across the grid
we tested, while accuracy stays within its interval.

All results are on a synthetic benchmark whose fraud parameters are modelling choices, not
observations; we state what a validation on real data would need to show.
\end{abstract}
